\documentclass[11pt]{article}
\usepackage[margin=0.88in]{geometry}
\usepackage{amsmath,amssymb,amsthm,mathtools}
\usepackage{lmodern,microtype,graphicx,xcolor,listings}
\usepackage[hidelinks]{hyperref}
\newtheorem{theorem}{Theorem}
\newtheorem{lemma}[theorem]{Lemma}
\newtheorem{proposition}[theorem]{Proposition}
\newcommand{\R}{\mathbb R}
\newcommand{\Q}{\mathbb Q}
\newcommand{\Z}{\mathbb Z}
\newcommand{\tr}{\operatorname{tr}}
\newcommand{\sgn}{\operatorname{sgn}}
\newcommand{\Ical}{\mathcal I_2}
\newcommand{\transpose}{^{\mathsf T}}
\DeclareMathOperator{\diag}{diag}
\setlength{\emergencystretch}{2em}
\lstset{language=Python,basicstyle=\ttfamily\scriptsize,
  columns=fullflexible,keepspaces=true,breaklines=true,
  showstringspaces=false,numbers=left,numberstyle=\tiny\color{gray},
  numbersep=8pt,frame=single,rulecolor=\color{black!20},
  xleftmargin=15pt,framexleftmargin=3pt,aboveskip=12pt,belowskip=12pt}
\title{A seven-variable piecewise quadratic function\protect\\
without a polynomial max--min representation}
\author{}
\date{}
\begin{document}
\maketitle
\vspace{-2em}
\begin{abstract}
An explicit rational linear subspace of pairs of matrices identifies a
continuous piecewise quadratic function on all of $\R^7$. The function has
six polynomial labels, including zero, and admits no finite expression
using maximum and minimum of real polynomials of arbitrary degrees.
The proof constructs, for every finite family of quadratic tests, two
points with the same test signs but different function signs. A finite
exact coefficient calculation and two applications of the real implicit
function theorem establish the required points. The coefficient calculation
is specified and its executable verifier is included. No minimality of
the dimension is asserted.
\end{abstract}

\section{Statement and the explicit function}
A function on $\R^n$ is piecewise polynomial here if a finite closed
semialgebraic cover carries polynomial formulas for it. A polynomial
lattice expression is a finite nonempty expression in real polynomials
using pointwise minimum and maximum. Distributivity puts such an expression
in finite maximum-of-minima form. Neither polynomial degrees nor the number
or real coefficients of the auxiliary polynomials are restricted.

\begin{theorem}\label{thm:main}
The function defined below is continuous and piecewise quadratic on all
of $\R^7$, with six distinct polynomial labels and a six-set closed
semialgebraic cover. It is not a polynomial lattice expression.
\end{theorem}

Let $A\in\R^{3\times3}$ and $B\in\R^{2\times3}$, with no symmetry
restriction on $A$. Number the ambient coordinates from zero, in row-major
order:
\[
w=(A_{00},A_{01},A_{02},A_{10},A_{11},A_{12},A_{20},A_{21},A_{22},
 B_{00},B_{01},B_{02},B_{10},B_{11},B_{12}).
\]
Reorder these as
\[
 P=(w_0,w_3,w_6,w_1,w_4,w_7,w_{11},w_{14}),\qquad
 N=(w_2,w_5,w_8,w_9,w_{12},w_{10},w_{13}).
\]
Define $T=M/630$, where
\begin{equation}\label{eq:M}
M=\begin{pmatrix}
-210&-210&-210&420&210&0&-210\\
326&-490&59&80&-185&3&80\\
-276&-210&66&-270&-150&-228&-270\\
266&-70&-91&-280&385&63&-280\\
-210&-210&-210&-210&210&0&420\\
-230&-140&370&-50&50&300&-50\\
36&0&279&90&225&153&90\\
-36&0&-279&-90&-225&477&-90
\end{pmatrix}.
\end{equation}
The vector space $W=\{P=TN\}$ is explicitly identified with $\R^7$ by
the \emph{unscaled} coordinates $z=N$. Thus every entry of $A(z),B(z)$
is a specified rational linear polynomial. Put
\[
 Q=A\transpose A-B\transpose B,\qquad H=\diag(1,1,-1).
\]
Direct rational multiplication gives $T\transpose T=I_7$. Since
$\tr(HQ)=|P|^2-|N|^2$, it follows that
\begin{equation}\label{eq:trace}
 Q_{00}+Q_{11}-Q_{22}=0\quad\hbox{on }W.
\end{equation}

Set $v_i=4e_i-\mathbf1$ for $0\leq i<3$, and $v_3=-\mathbf1$, where
$\mathbf1=(1,1,1)\transpose$. Define
\[
 q_{ij}=-v_i\transpose Qv_j\quad(i<j),\qquad
 E=\{01,02,03,12,13\},\qquad h=\min_{e\in E}q_e.
\]
The function of the theorem is
\begin{equation}\label{eq:f}
 f(z)=\begin{cases}h(z),&h(z)>0\text{ and }\det A(z)>0,\\
                    0,&\text{otherwise}.
       \end{cases}
\end{equation}
This definition applies to every $z\in\R^7$.

\subsection{Positivity, continuity, and the closed cover}
Let $u_i=(e_i+\mathbf1)/4$ for $i<3$ and $u_3=0$.
The matrix $(v_i\transpose Qv_j)_{i,j=0}^3$ has zero row sums,
off-diagonal entries $-q_{ij}$, and $\sum_i u_iv_i\transpose=I_3$.
Consequently
\begin{equation}\label{eq:simplex}
 Q=\sum_{i<j}q_{ij}(u_i-u_j)(u_i-u_j)\transpose.
\end{equation}
Taking $H$-traces and using \eqref{eq:trace} gives
\begin{equation}\label{eq:labelrelation}
 q_{23}=q_{01}+2q_{03}+2q_{13}.
\end{equation}
If $h>0$, all six coefficients in \eqref{eq:simplex} are positive.
The vectors $u_0,u_1,u_2$ span $\R^3$, so $Q$ is positive definite.
Then $A$ is invertible: $Av=0$ would imply
$v\transpose Qv=-|Bv|^2\leq0$. Thus $\det A$ has locally constant
sign wherever $h>0$. The function is continuous there and on $\{h<0\}$;
at $h=0$ continuity follows from $0\leq f\leq\max(h,0)$.
Also
\begin{equation}\label{eq:homog}
 f(rz)=r^2f(z)\qquad(r>0).
\end{equation}

For the asserted cover, take
\[
 F_0=\{\det A\leq0\}\ \cup\!\bigcup_{e\in E}\{q_e\leq0\}
\]
with label zero and, for each $e\in E$, take
\[
 F_e=\{\det A\geq0,\ q_b\geq0\ (b\in E),\ q_e\leq q_b\ (b\in E)\}
\]
with label $q_e$. These sets are closed and semialgebraic and cover $W$.
If $\det A=0$ and all five $q_e$ are nonnegative, their minimum is zero
by the preceding positivity argument. This checks the formula on every
boundary overlap. Distinctness of the five nonzero labels will also follow
from the rank calculation below.

\section{A real curve and its finite coefficient calculation}
It is convenient to specify the linear constraint map by
\[
 (Lw)_i=s_i\bigl(P_i-(TN)_i\bigr),\qquad
 s=(3,630,105,90,3,63,70,70).
\]
Its entries are integers and $\ker L=W$. The full matrix appears in
the verifier in Appendix~\ref{app:verifier}.
The base point is
\begin{equation}\label{eq:base}
 A_0=\diag(1,-1,0),\qquad
 B_0=\begin{pmatrix}1&0&0\\0&-1&0\end{pmatrix}.
\end{equation}
It lies in $W$, and its $N$-coordinates are $(0,0,0,1,0,0,-1)$.

For $k=(k_0,k_1,k_2)$ define
\[
 d=1+|k|^2,\qquad
 R_k=(1-|k|^2)I_3+2kk\transpose+2[k]_\times,
 \qquad [k]_\times v=k\times v.
\]
Direct multiplication gives $R_k\transpose R_k=d^2I_3$ and
$\det R_k=d^3$. For the latter identity one can also square the
determinant and use its nonvanishing and its positive sign at $k=0$.
Thus $O_k=R_k/d\in SO(3)$. For $C\in\R^{2\times2}$ set
\begin{equation}\label{eq:chart}
 V=\begin{pmatrix}1&0&t_0\\0&1&t_1\end{pmatrix},\qquad
 \Phi_A=R_k[:,0\!:\!2]CV,\qquad \Phi_B=dCV.
\end{equation}
Here $R_k[:,0\!:\!2]$ denotes the first two columns. The chart is
polynomial in nine parameters and satisfies
$\Phi_A\transpose\Phi_A=\Phi_B\transpose\Phi_B$ identically.
If $C$ is invertible, both matrices have rank two, and
$\Phi_A=O_k[\Phi_B;0]$.

Fix $C_{00}=1$, put $C_{11}=-1+t$, and regard
$u=(k_0,k_1,k_2,t_0,t_1,C_{01},C_{10})$ as seven unknowns.
Let $F(t,u)$ consist of rows 1 through 7 of $L\Phi$, with row numbering
starting at zero. Then $F(0,0)=0$ and
\begin{equation}\label{eq:J}
 J=D_uF(0,0)=\begin{pmatrix}
0&0&1260&-326&-490&-3&815\\
0&-210&0&46&-35&38&25\\
0&0&180&-38&-10&81&-55\\
0&0&0&1&-1&0&-1\\
-126&0&0&23&-14&-30&-5\\
0&0&0&66&0&-17&-25\\
0&0&0&4&-70&-53&25
\end{pmatrix},
\end{equation}
with
\begin{equation}\label{eq:D}
 D=\det J=16383079440000\ne0.
\end{equation}
The real implicit function theorem supplies a smooth $u(t)$ near zero
with $u(0)=0$ and $F(t,u(t))=0$.

The omitted equation must be recovered. The active seven equations imply
$P=TN+\ell e_0$. On the Gram-zero chart,
\begin{equation}\label{eq:omitted}
 0=\tr(HQ)=\ell\bigl(2(TN)_0+\ell\bigr).
\end{equation}
The second factor is 2 at \eqref{eq:base} and remains nonzero near it.
Hence $\ell=0$. Write $b(t)=\Phi(t,u(t))$, as a curve in $W$ or in its
seven coordinates. It has Gram matrix zero and rank-two matrices for
all sufficiently small $t$.

\subsection{The specified finite calculation}
Over $\Q$, recursively define
$u_{22}(t)=\sum_{j=1}^{22}u_jt^j$ by
\begin{equation}\label{eq:recursion}
 u_j=-J^{-1}[t^j]\,F\!\left(t,\sum_{a<j}u_at^a\right).
\end{equation}
An unknown coefficient of order $j$ enters the order-$j$ equation only
through $J$. Thus $F(t,u_{22}(t))$ is divisible by $t^{23}$, and all
coefficients of $u_{22}$ belong to $\Z[1/D]$.
Substitute this polynomial in \eqref{eq:chart}. Let $C_1$ be the
$23\times7$ matrix of the $N$-coordinate coefficients of orders 0 through
22, and let $C_2$ be the $23\times28$ matrix for their products $N_iN_j$,
ordered lexicographically for $0\leq i\leq j<7$.
Let $G_2$ be the $6\times28$ matrix of the Gram entries, in the order
$00,01,02,11,12,22$, as quadratic polynomials in $N$.

The following are determinants of fixed minors, with zero-based indices,
reduced modulo the prime 1009:
\begin{align}
 \det C_1[0\!:\!7,\,0\!:\!7]&=693,\label{eq:c1}\\
 \det C_2[0\!:\!23,\,0\!:\!23]&=349,\label{eq:c2}\\
 \det G_2[0\!:\!5,\,0\!:\!5]&=256.\label{eq:g2}
\end{align}
Each interval includes its left endpoint and excludes its right endpoint.
Appendix~\ref{app:verifier} implements the complete calculation by
truncated polynomial multiplication and elimination. In particular, no
sampling of a numerical real curve is used. The prime does not divide $D$
or the denominators of $T$. Therefore reduction is defined on all
coefficients in the recursion and the three displayed rational minors.
Their nonzero residues prove nonvanishing over $\Q$ and over $\R$.

For clarity, the finite polynomial approximation determines the required
coefficients of the \emph{actual real curve}, without presupposing convergence
of a formal power series. Set $\widetilde F=J^{-1}F$. In a sufficiently
small convex neighborhood,
$\|D_u\widetilde F-I\|<1/2$. Integration along the segment from
$u(t)$ to $u_{22}(t)$ gives
\begin{equation}\label{eq:error}
 \|u_{22}(t)-u(t)\|
 \leq2\|\widetilde F(t,u_{22}(t))\|=O(|t|^{23}).
\end{equation}
Both endpoints tend to zero, so the segment lies in this neighborhood
for small $t$. Polynomial substitution in the chart, or into any fixed
quadratic test, preserves this error order.

Let $\Ical$ denote the real span of the six Gram quadratics on $W$.
Equation~\eqref{eq:trace} and \eqref{eq:g2} give $\dim\Ical=5$.
Every member vanishes on $b(t)$, hence belongs to $\ker C_2$ by
\eqref{eq:error}. Equation~\eqref{eq:c2} and $28-5=23$ give
\begin{equation}\label{eq:kernels}
 \ker C_1=0,\qquad \ker C_2=\Ical.
\end{equation}
The six simplex labels are an invertible linear change of the six entries
of a symmetric matrix, by \eqref{eq:simplex}. Their unique restricted
relation is \eqref{eq:labelrelation}, whose coefficient of $q_{23}$ is
nonzero. Thus the five labels used in $f$ are linearly independent.

\section{Avoiding any finite family of quadratic tests}
Fix an arbitrary finite family $\mathcal T$ of real polynomials on $W$
of degree at most two. For $p=c+\lambda+q\notin\Ical$, with homogeneous
components of degrees zero, one, and two, the order-$j$ coefficient of
$p(\rho b(t))$ through order 22 is
\begin{equation}\label{eq:rho}
 \delta_{j0}c+\rho[\lambda(b(t))]_j+\rho^2[q(b(t))]_j.
\end{equation}
At least one of these is a nonzero polynomial in $\rho$. Otherwise
comparison of its degrees, followed by \eqref{eq:kernels}, would give
$c=0$, $\lambda=0$, and $q\in\Ical$.
For each of the finitely many $p\notin\Ical$, choose one such nonzero
coefficient polynomial. Choose $\rho>0$ close to 1 outside their finite
sets of roots. For this fixed $\rho$, \eqref{eq:error} and the first
nonzero coefficient show that every one of these tests is nonzero at
$\rho b(t)$ for all sufficiently small positive $t$. Choose a common
such $t$ and put
\begin{equation}\label{eq:bscaled}
 b=\rho b(t).
\end{equation}
Both choices may be made in any prescribed neighborhoods of 1 and 0.
Thus the local rank and derivative conditions needed below hold at this
base. The tests in $\Ical$ vanish there; all remaining tests do not.

\section{Two paths with the same exact Gram matrix}
Fix the matrices
\begin{equation}\label{eq:target}
 R=\begin{pmatrix}1&1/4&0\\1/4&1&0\\0&0&2\end{pmatrix},\qquad
 S=\begin{pmatrix}1&1/4&0\\0&\sqrt{15}/4&0\\0&0&\sqrt2\end{pmatrix}.
\end{equation}
They satisfy $S\transpose S=R$, $\det S>0$, $R\succ0$ and
$\tr(HR)=0$. The six simplex labels of $R$, in order $01,02,03,12,13,23$,
are
\begin{equation}\label{eq:targetlabels}
 (3/2,17/2,1/2,17/2,1/2,7/2).
\end{equation}

Keep the chosen $t,\rho$ fixed, but temporarily allow $u$ to vary near
$u(t)$. Set $B=\rho\Phi_B$ and $O=O_k$, so the moving base matrix
$A=\rho\Phi_A$ is $O[B;0]$, irrespective of whether the linear
constraints hold. Define
\[
 B'=BS^{-1},\quad G=B'(B')\transpose,\quad
 n=\frac{\operatorname{row}_0(B')\times\operatorname{row}_1(B')}
 {|\operatorname{row}_0(B')\times\operatorname{row}_1(B')|}.
\]
Nearby, $B'$ has row rank two, $G\succ0$, and these functions are smooth.
For small signed $\varepsilon$, put $M_\varepsilon=I_2-\varepsilon^2G^{-1}$
and
\begin{equation}\label{eq:sqrt}
 U=\frac{M_\varepsilon+\sqrt{\det M_\varepsilon}\,I_2}
 {\sqrt{\tr M_\varepsilon+2\sqrt{\det M_\varepsilon}}}.
\end{equation}
All radicands are strictly positive nearby. The two-by-two
Cayley--Hamilton identity gives $U^2=M_\varepsilon$; $U$ is symmetric
and equals $I_2$ at $\varepsilon=0$. Define
\begin{equation}\label{eq:rawpaths}
 \widehat A_\varepsilon
 =O\begin{pmatrix}B'\\\varepsilon n\transpose\end{pmatrix}S,
 \qquad \widehat B_\varepsilon=UB'S.
\end{equation}
At zero these equal the moving scaled base chart.

The orthogonal projection onto the row space of $B'$ is both
$(B')\transpose G^{-1}B'$ and $I_3-nn\transpose$. Therefore, for all
nearby parameters, not merely those satisfying the constraints,
\begin{align}
 \widehat A_\varepsilon\transpose\widehat A_\varepsilon
 -\widehat B_\varepsilon\transpose\widehat B_\varepsilon
 &=\varepsilon^2S\transpose
       (nn\transpose+(B')\transpose G^{-1}B')S
   =\varepsilon^2R,\label{eq:exactgram}\\
 \det\widehat A_\varepsilon
 &=\varepsilon\,|\operatorname{row}_0(B')\times
                     \operatorname{row}_1(B')|\,\det S.
 \label{eq:exactdet}
\end{align}

Impose rows 1 through 7 of $L$ on \eqref{eq:rawpaths} and solve for $u$
as a function of $\varepsilon$. At zero the equation holds at $u(t)$;
its $u$-derivative is $\rho D_uF(t,u(t))$. It is invertible for
sufficiently small $t$ by \eqref{eq:D} and continuity. The real implicit
function theorem yields a smooth solution for both signs of small
$\varepsilon$. By \eqref{eq:exactgram}, the perturbed $H$-trace is zero.
The active equations again give $P=TN+\ell e_0$, so
\eqref{eq:omitted} still holds. Its second factor is close to $2\rho>0$;
the omitted row therefore holds as well.

We have obtained a path $x(\varepsilon)$ in the \emph{same fixed space}
$W$, with $x(0)=b$ and
\begin{equation}\label{eq:finalpaths}
 Q(x(\varepsilon))=\varepsilon^2R,\qquad
 \sgn\det A(x(\varepsilon))=\sgn\varepsilon\quad(\varepsilon\ne0).
\end{equation}
In particular, for sufficiently small $\varepsilon>0$,
\begin{equation}\label{eq:values}
 f(x(\varepsilon))=\varepsilon^2/2>0,\qquad f(x(-\varepsilon))=0.
\end{equation}
For every $p\in\mathcal T\setminus\Ical$, $p(b)\ne0$ and continuity
gives the same nonzero sign at both points. For every $p\in\mathcal T
\cap\Ical$, the exact Gram identity gives equal values at the two
points, including when that common value is zero. Finiteness permits
one positive $\varepsilon$ working for all tests. We have proved:

\begin{proposition}\label{prop:signs}
For every finite family of real polynomials of degree at most two on $W$,
there are two points at which all members of the family have identical
signs, but $f$ has different signs.
\end{proposition}
The order of choices was: fixed $W,f$; arbitrary finite family;
$\rho$; a positive $t$ and its base; corrected paths; a positive
$\varepsilon$. Neither $W$ nor $f$ depends on the tests.

\section{Removing every degree bound on a representation}
\begin{lemma}\label{lem:radial}
If a positively homogeneous function of degree two is a polynomial lattice
expression, its sign is determined by the signs of finitely many real
polynomials of degree at most two.
\end{lemma}
\begin{proof}
Write $g=\mathcal L(p_1,\ldots,p_m)$ for a finite min--max expression,
and decompose each polynomial as
$p_i=c_i+\lambda_i+q_i+$ terms of degree at least three.
For fixed $x$, as $r\to0^+$, the limit of $r^{-2}p_i(rx)$ in the
extended real line is
\[
 \begin{cases}
  \sgn(c_i)\,\infty,&c_i\ne0,\\
  \sgn(\lambda_i(x))\,\infty,&c_i=0,\ \lambda_i(x)\ne0,\\
  q_i(x),&c_i=0,\ \lambda_i(x)=0.
 \end{cases}
\]
Positive scaling commutes with minimum and maximum. Homogeneity gives
\[
 g(x)=\mathcal L(r^{-2}p_1(rx),\ldots,r^{-2}p_m(rx))\qquad(r>0).
\]
The two lattice operations are continuous on the extended real line,
so this equality passes to their finitely many limits. After this passage,
apply the sign map to the result: the sign map commutes with minimum and
maximum in the ordered sets $[-\infty,\infty]$ and $\{-1,0,1\}$.
The sign of $g(x)$ is therefore determined by the signs of the finite
family $\{\lambda_i,q_i:1\leq i\leq m\}$ and the fixed constants $c_i$.
No continuity of the sign map is asserted or used.
\end{proof}

Applying Lemma~\ref{lem:radial} to \eqref{eq:homog} contradicts
Proposition~\ref{prop:signs}. This proves Theorem~\ref{thm:main}, including
the absence of any bound on the number, real coefficients, or degrees of
the auxiliary polynomials. The argument proves a finite quadratic
\emph{sign} condition; it does not assert that a lattice representation
could be replaced by one whose leaves all have degree at most two.

\appendix
\section{Executable finite certificate}\label{app:verifier}
The following Python 3 program uses NumPy and the standard library.
It performs no random generation, numerical rank estimates, or file writes.
The isometry check uses exact fractions. The Jacobian determinant uses
integer Bareiss elimination. All truncated polynomial operations and
subsequent elimination are modulo 1009; each multiplication is reduced
before another multiplication. At order 22, the displayed dimensions and
bounded residues keep all NumPy integer intermediates below $2^{63}$.
The first derivatives of the chart at its base have entries in
$\{0,\pm1,\pm2\}$, so the centered modular recovery used for that
integer Jacobian is exact. The characteristic-zero lifting and the
comparison with the real curve are proved in the body of the paper.

The expected rank and determinant triples are $(7,693)$, $(23,349)$,
and $(5,256)$. The output also gives the selected row and column indices;
they are respectively the three fixed minors
\eqref{eq:c1}--\eqref{eq:g2}. A failed check raises an exception, including
when Python is invoked with optimization enabled.

\lstinputlisting{verify_certificate7.py}
\end{document}
